\chapter{实作：从粗混到成品}
\label{chap:rough-to-finished}

这一章不再按器件讲，而是把一首歌从头走到尾。示例是一个 3 分 50 秒的工程：鼓、贝斯、两把吉他、主唱、和声、弦乐，段落为前奏、主歌一、副歌一、主歌二、副歌二、桥段、终副歌、尾奏。下面的数值是经验起点，不是夜半内置预设；每一步都写清为什么动、动到哪、听到什么才算过。

\begin{center}
\begin{tikzpicture}[
  box/.style={draw, rounded corners=2pt, minimum width=27mm, minimum height=8mm, align=center, font=\scriptsize},
  arrow/.style={->, thick}
]
  \node[box] (prep) at (0,0) {1 准备\\与参考};
  \node[box] (rough) at (3.4,0) {2 粗平衡\\推子/声像};
  \node[box] (gain) at (6.8,0) {3 增益结构\\峰值与余量};
  \node[box] (eq) at (10.2,0) {4 EQ\\低架先行};
  \node[box] (dyn) at (10.2,-1.8) {5 动态/空间\\压缩与混响};
  \node[box] (auto) at (6.8,-1.8) {6 自动化\\段落曲线};
  \node[box] (bus) at (3.4,-1.8) {7 总线\\共享处理};
  \node[box] (out) at (0,-1.8) {8 导出\\检查与实测};
  \draw[arrow] (prep) -- (rough);
  \draw[arrow] (rough) -- (gain);
  \draw[arrow] (gain) -- (eq);
  \draw[arrow] (eq) -- (dyn);
  \draw[arrow] (dyn) -- (auto);
  \draw[arrow] (auto) -- (bus);
  \draw[arrow] (bus) -- (out);
\end{tikzpicture}
\end{center}

八步不是八个孤立的按钮，而是一条回退链：任何一步听不对，先回到上一步，不要在总线上硬补。下面的段落表给出示例曲的检查点；真实作品可以换时间，但每个段落都要有“这一段要比上一段多什么或少什么”的一句话目标。

\begin{table}[htbp]
\centering
\small
\begin{tabular}{@{}p{0.16\linewidth}p{0.15\linewidth}p{0.34\linewidth}p{0.27\linewidth}@{}}
\hline
段落 & 时间 & 这一段的目标 & 主要动作 \\
\hline
前奏 & 0:00--0:35 & 建立脉冲，留足余量 & 粗平衡，不加总线处理 \\
主歌一 & 0:35--1:10 & 主唱字头清楚，鼓与贝斯不抢 & 推子、低架、主唱压缩 \\
副歌一 & 1:10--1:45 & 比主歌高一档，但不靠推子硬顶 & 主唱自动化、空间、总线胶水 \\
主歌二 & 1:45--2:20 & 换纹理，避免和主歌一重复 & 吉他延迟或换声部、主唱动态 \\
副歌二 & 2:20--2:55 & 比副歌一更大，层次更多 & 和声、双轨、辅助返回 \\
桥段 & 2:55--3:15 & 突然收小，暴露主唱 & 鼓组退让、主唱推起、混响收短 \\
终副歌 & 3:15--3:35 & 全曲最大，但仍不糊 & 全轨进入、总线限幅兜底 \\
尾奏 & 3:35--3:50 & 尾巴干净，不突然断 & 混响尾、音量收尾 \\
\hline
\end{tabular}
\caption{示例曲的段落目标}
\label{tab:walk-sections}
\end{table}

\section{准备与参考曲}

先把工程存一次，再确认音频设置里的采样率；导出对话框的采样率默认是 \texttt{auto}，跟随工程，不要在导出时才改采样率。参考曲走“文件 → 导入音频文件...”，夜半的 GUI 导入只接受 WAV、AIF、AIFF；其他格式要走 MCP 的导入路径，不能把“文件选择器没反应”当成工程坏了。

导入后，选中一条音频轨，走“轨道 → 将最新导入音频摆放到当前音频轨”。参考曲没有独立的“参考轨”类型，也没有自动响度匹配：它就是一条普通音频轨，仍然经过主总线。A/B 时用独奏/静音切换，并把这条轨的设备旁通；不要说自己有一条“绕过主总线的参考总线”，夜半没有这个入口。

参考轨的正确用法是定“方向”，不是定“数值”：听它的低频密度、主唱在混音里的位置、副歌的抬升幅度，再回到自己的工程做相对判断。不要拿参考曲的 LUFS 直接当自己的目标，也不要把参考轨留在最终导出里。

\section{粗平衡与增益结构}

先独奏鼓、贝斯、主唱三组，把别的轨静音。推子从 $-12$ dB 附近开始，主总线保持在 $0$ dB；先把这三组的峰值关系摆对，再放吉他和弦乐。不要一上来把所有推子推到 $0$ dB：推子位置不是“标准答案”，只是电平控制。

\begin{table}[htbp]
\centering
\small
\begin{tabular}{@{}p{0.24\linewidth}p{0.22\linewidth}p{0.46\linewidth}@{}}
\hline
控制 & 夜半范围 & 用法 \\
\hline
轨道推子 & $-60$ 至 $+6$ dB & 粗平衡从 $-12$ dB 起；不要用它修源素材的削波 \\
声像 & $-1$ 至 $+1$ & 默认 $0$；节拍骨架留在中间，吉他与和声少量偏置 \\
发送增益 & $-60$ 至 $+12$ dB & 绝对 dB；$0$ dB 是单位增益，不是“百分之多少” \\
主总线推子 & $-60$ 至 $+6$ dB & 混音期保持 $0$；推高会让主插入链与导出限制器收到更大电平 \\
轨道电平表 & 峰值/RMS 均为 dBFS & 经验起点：底鼓约 $-6$、贝斯约 $-9$、主唱约 $-6$、吉他约 $-12$；不是自动目标 \\
主总线电平表 & 峰值/RMS 为 dBFS，LUFS 为 RMS 派生 & 混音期把主峰值留在 $-6$ 至 $-3$ dBFS，给总线处理留余量 \\
\hline
\end{tabular}
\caption{粗平衡涉及的控制与经验起点}
\label{tab:walk-controls}
\end{table}

增益结构不是“把每一轨都推到一样响”。夜半的通道条有独立的输入增益与输出增益，范围都是 $-60$ 至 $+24$ dB；今天这两个值没有 GUI 滑杆，只能在 MCP 或工程 JSON 里写。界面上的推子是另一层电平控制。推子拉低不等于把源素材的削波修好，推子推高也不等于给设备链更多余量；先看轨道峰值表，再决定是换素材、改通道条增益，还是改推子。

经验上，24-bit 工程不必把每条轨都压到 $0$ dBFS。让底鼓、贝斯、主唱的峰值分别落在 $-9$ 至 $-3$ dBFS 的区间，主总线在进入总线处理前留 $-6$ dB 左右；这样后面加 EQ 的提升、压缩的补偿增益、混响的湿声，都不容易一推就顶到天花板。以上是起点，不是夜半的自动归一化，夜半也不会替你判断“够不够响”。

粗平衡的听感判据：在主歌一，不看字幕能听清主唱的每个词；底鼓和贝斯同时出现时，贝斯的低端不把底鼓的冲击糊掉；副歌一第一次出现时，整体比主歌明显抬一档，但主总线峰值没有贴到 $0$。如果做不到，先回推子与声像，不要先加 EQ。

\screenshot{walkthrough-console}{混音控制台至少六条轨，推子、声像、轨道电平与主总线电平同屏}

\section{EQ、动态与空间}

EQ 的第一原则是“先减后加”。夜半的 EQ 器件有三段：低架、中频峰、高架；参数范围是低架增益 $-18$ 至 $+18$ dB、低架转折 $20$ 至 $1000$ Hz，中频增益 $-18$ 至 $+18$ dB、中心 $200$ 至 $8000$ Hz、$Q$ $0.3$ 至 $6$，高架增益 $-18$ 至 $+18$ dB、转折 $1000$ 至 $16000$ Hz。默认全部为 $0$ dB，低架 $200$ Hz、中频 $1000$ Hz/$Q=0.9$、高架 $4000$ Hz。

但界面今天只暴露低架增益一个滑块：装上 EQ 后，设备机架里能拖的是 \texttt{eq\_low\_gain}，低架转折仍是默认的 $200$ Hz；中频、高架、频率与 $Q$ 都没有 GUI 控件，只能在 MCP 或工程 JSON 里写。不要把“装 EQ”理解成“四段 EQ 都能在界面上拧”；界面没有四个旋钮，也没有频率读数。

因此 GUI 内的 EQ 走查要收敛：先用低架增益做 $1$ 至 $3$ dB 的微调，判断“加了之后低频更清楚，还是只是更浑”。

\screenshot{walkthrough-eq-low}{设备机架 EQ 面板只显示低架增益滑块，频率与中高架不可见}

如果目标是 $80$ 至 $120$ Hz 的贝斯低端，GUI 的 $200$ Hz 低架不是同一件事；要么走 MCP 把 \texttt{eq\_low\_freq} 改到目标，要么承认这一版 GUI 做不到。中频刺耳、齿音、鼻音这些窄问题，优先用中频峰的 $1$ 至 $2$ dB 衰减，$Q$ 从 $0.9$ 开始，不要用宽低架去压。

\begin{table}[htbp]
\centering
\small
\begin{tabular}{@{}p{0.13\linewidth}p{0.50\linewidth}p{0.16\linewidth}p{0.19\linewidth}@{}}
\hline
器件 & 关键参数与值域 & 缺省 & 界面可达 \\
\hline
EQ & 低架增益 $-18$ 至 $+18$ dB、低架转折 $20$ 至 $1000$ Hz；中频增益 $-18$ 至 $+18$ dB、中心 $200$ 至 $8000$ Hz、$Q$ $0.3$ 至 $6$；高架增益 $-18$ 至 $+18$ dB、转折 $1000$ 至 $16000$ Hz & 增益全 $0$；低架 $200$ Hz；中频 $1000$ Hz、$Q=0.9$；高架 $4000$ Hz & 只有低架增益滑块；频率、中频、高架走 MCP/JSON \\
压缩器 & 阈值 $-120$ 至 $+24$ dBFS、比率 $1$ 至 $1000$、膝宽 $0$ 至 $24$ dB、检波 $0.0001$ 至 $10$ s、起振 $0.0001$ 至 $10$ s、释放 $0.0001$ 至 $10$ s、补偿 $-24$ 至 $+24$ dB & 阈值 $-12$、比率 $4$、膝 $6$ dB、检波 $5$ ms、起振 $10$ ms、释放 $100$ ms、补偿 $0$ & 轨道菜单可插；设备机架只能循环 $2/4/8/20$ 比率；其余参数走 MCP/JSON \\
\hline
\end{tabular}
\caption{EQ 与压缩器的参数核对}
\label{tab:walk-eq-comp}
\end{table}

压缩器装上去的默认值是阈值 $-12$ dBFS、比率 $4:1$、膝宽 $6$ dB、检波 $5$ ms、起振 $10$ ms、释放 $100$ ms、补偿 $0$ dB。GUI 能做的只有两件事：从轨道菜单插一台，或者在设备机架里点压缩器行循环 $2/4/8/20$ 的比率；没有阈值、起振、释放、补偿的滑杆。设备机架里那行“压缩器: 阈值 ... 比率 ... GR 0.0 dB”的 $GR$ 是固定文案，不是实时增益衰减表；不能用它判断“压了多少”。要精确设阈值与时间常数，走 MCP；要听，就用旁通 A/B，判断压缩后瞬态是否还在、音节是否变糊。

压缩的听感判据不是“GR 数字越大越好”。

\screenshot{walkthrough-compressor}{压缩器行显示阈值、比率与 GR 文案；GR 文案不是实时表}

主唱压缩后，辅音仍然清楚、长音不抽吸、句子之间的呼吸没有被抹平；贝斯压缩后，低频不“泵”，但峰不再顶到主总线；总线压缩后，整首更粘连，而不是所有段落被压成同一个响度。没有真实 GR 表时，不要硬追某个数字。

空间处理先分插入式与发送式。插入式混响装在轨上，干声和湿声一起过这只混响；发送式混响要把混响装在辅助返回轨上，原轨只送一份信号过去。夜半的混响默认是 Room 类型、\texttt{reverb\_size} 为 $0.45$、\texttt{reverb\_damp} 为 $0.35$、\texttt{reverb\_wet} 为 $0.25$、预延迟 $12$ ms；\texttt{reverb\_wet} 同时决定干湿比例，$0.25$ 表示约四分之一的湿声、四分之三的干声。

这就是共享混响的坑：把默认 \texttt{reverb\_wet}（$0.25$）的混响直接装在辅助返回上，辅助返回会把 $75\%$ 的干声再送进主总线，和原轨的干声叠成双干声。共享混响要把 \texttt{reverb\_wet} 推到接近 $1.0$，让干路接近 $0$；这个参数今天没有 GUI 滑杆，只能走 MCP/JSON。GUI 的“装混响”适合轨上插入式混响，不适合直接当共享返回用。

\begin{table}[htbp]
\centering
\small
\begin{tabular}{@{}p{0.13\linewidth}p{0.55\linewidth}p{0.15\linewidth}p{0.17\linewidth}@{}}
\hline
器件 & 关键参数与值域 & 缺省 & 界面可达 \\
\hline
混响 & 大小 $0$ 至 $1$、类型 $0$ 至 $3$（Room/Plate/Hall/Spring）、T60 $0$ 至 $10$ s、阻尼 $0$ 至 $1$、湿声比例 $0$ 至 $1$、预延迟 $0$ 至 $0.1$ s & Room、大小 $0.45$、阻尼 $0.35$、湿声 $0.25$、预延迟 $12$ ms & 可插入、可旁通；参数走 MCP/JSON \\
延迟 & 时间 $0.001$ 至 $2$ s、反馈 $0$ 至 $0.95$、湿声 $0$ 至 $1$、阻尼 $0$ 至 $1$、乒乓 $0/1$、同步 $0/1$、分子/分母 $1$ 至 $16$、附点 $0/1$ & 时间 $0.3$ s、反馈 $0.3$、湿声 $0$、乒乓 $1$、同步关 & 没有装延迟的 GUI 入口；已在工程里时可旁通，参数走 MCP/JSON \\
限制器 & 阈值 $-30$ 至 $0$ dB、天花板 $-6$ 至 $0$ dB、释放 $10$ 至 $1000$ ms、前瞻 $1$ 至 $192$ 帧、真峰值 $0/1$ & 阈值 $-1$、天花板 $-0.3$、释放 $80$ ms、前瞻 $48$、真峰值开 & 没有装限制器的 GUI 入口；参数走 MCP/JSON \\
\hline
\end{tabular}
\caption{空间、延迟与限制器的参数核对}
\label{tab:walk-space-limiter}
\end{table}

主唱与吉他的空间不要同时堆在一条轨上：主唱用短混响，吉他可以用延迟；两者都湿，副歌会糊。单声道兼容检查按 F4 打开监听控制室，切 MONO 听低频与和声是否消失，切 SIDE 只听左右差值；这两个开关只改监听，不改导出。若 MONO 下贝斯或主唱塌掉，回到声像与宽度，不要用总线 EQ 硬补。

\section{自动化}

自动化的作用是把“段落之间的差异”写死，而不是把每一拍都画一遍。夜半的自动化目标是音量、声像、发送增益、设备参数、宏与 MIDI CC；今天 GUI 只能点击已经存在的泳道来加/改点，不能新建或删除泳道、不能拖点、不能改曲线、不能切读/写模式。没有泳道时，先在 MCP 或工程 JSON 里建；不要对着空白编曲区反复点，界面不会凭空长出一条 lane。

音量泳道的标称纵轴是 $-60$ 至 $+12$ dB，界面推子本身是 $-60$ 至 $+6$ dB；这两个值域不是同一个东西。常规推子范围内的段落抬升用 $+1$ 至 $+3$ dB，不要把自动化写成推子到不了的 $+12$ 再来问“为什么声音不对”。写点前先把轨道静态推子定在目标附近，再看自动化点是不是绝对 dB；读取开关默认开，关掉后播放与导出回退到静态值。

\begin{table}[htbp]
\centering
\small
\begin{tabular}{@{}p{0.23\linewidth}p{0.17\linewidth}p{0.18\linewidth}p{0.40\linewidth}@{}}
\hline
目标 & 段落 & 值 & 理由 \\
\hline
主唱音量 & 副歌一 & $+1.5$ dB & 副歌抬一档，但不靠整轨推子硬顶 \\
主唱音量 & 桥段 & $-3$ dB & 收小，让吉他/弦乐暴露 \\
主唱音量 & 终副歌 & $+2$ dB & 全曲最大，但仍留总线余量 \\
吉他音量 & 主歌二 & $-2$ dB & 换纹理，不和主歌一抢 \\
混响发送 & 副歌一 & $-6$ dB & 只加尾巴，不把干声再送一份 \\
混响发送 & 桥段 & $-12$ dB & 收短空间，突出主唱 \\
\hline
\end{tabular}
\caption{示例曲的自动化点}
\label{tab:walk-automation}
\end{table}

自动化点的写法：每个段落 $2$ 至 $4$ 个点足够，点在段落边界与情绪转折处；不要在每小节都写点，否则曲线会变成噪声。点与点之间按曲线插值，端点精确；界面点击既有点的 $8$ 像素范围会改原点，不会新建重复点。听感判据是“听不出自动化在动”：主唱从主歌到副歌的抬升应该像演唱自然变强，而不是像推子被拉了一下。

\screenshot{walkthrough-automation}{编曲区音量泳道有段落点，主唱轨与副歌边界可见}

\section{总线与导出}

总线阶段先分清两条腿：辅助返回轨用来做“这一组共享”的处理，主总线用来做“整首最后一道”。新建辅助返回轨走“轨道 → + 新建辅助总线轨”，把组内轨的发送边指向它；发送边只有增益，没有 pre/post 开关，别按其他 DAW 的发送点习惯假设。

主总线的 GUI 设备机架只给装 EQ 与混响两只按钮，没有主总线压缩器按钮；要做“总线胶水”，选一条辅助返回轨，在轨道菜单里给它插 Compressor，再把鼓组或全组发送进去。主总线自己的 EQ 低架滑块同样只改低架增益，频率与中高架走 MCP。主总线推子保持 $0$ dB；推高会让插入链与导出限制器收到更大电平，不要用推子给限制器“做补偿”。

导出阶段有两只限制器，必须分开说。设备链里的 Limiter 是普通插入器件：阈值默认 $-1$ dBFS、天花板 $-0.3$ dBFS、释放 $80$ ms、前瞻 $48$ 帧、真峰值检测开；今天 GUI 没有装它的按钮，参数也只能走 MCP/JSON。导出对话框里的“母线限制器”是另一只固定参数限制器：阈值 $0.9$（约 $-0.92$ dBFS）、天花板 $0.95$（约 $-0.45$ dBFS）、前瞻 $33$ 帧、真峰值检测关；默认关，每次导出一个布尔量，不写进工程，也不进撤销栈。

\begin{table}[htbp]
\centering
\small
\begin{tabular}{@{}p{0.18\linewidth}p{0.29\linewidth}p{0.45\linewidth}@{}}
\hline
控件 & 默认 & 说明 \\
\hline
输出路径 & 空 = 宿主缺省 & 可手敲或“浏览...”选择；只填框，不自动导出 \\
采样率 & \texttt{auto} & 跟随工程；只有确实要改时才填 Hz \\
位深 & \texttt{24} & 16/24 整数走 TPDF 抖动；32f 浮点透传、不抖动 \\
时长 & 工程可闻末端 & 清空会回退 8 秒缺省；不是“选中片段” \\
母线限制器 & 关 & 仅本次导出；固定参数，不能当响度归一化 \\
\hline
\end{tabular}
\caption{WAV 导出对话框的真实控件}
\label{tab:walk-export-controls}
\end{table}

\screenshot{walkthrough-export}{WAV 导出对话框的五个控件与导出后实测行同屏}

导出后对话框会写一行“实测 · 真峰值 ... dBFS · 积分响度 ... LUFS · LRA ... LU”。三个值都来自同一次母带测量；测不出就写“未测出”，不要把它当成 $0$。真峰值用 $16$ 倍过采样，积分响度是带门限的积分，LRA 是短时响度的范围；它们不是控制室主表那个“平滑 RMS 减 $0.691$ dB”的 LUFS。若要在流媒体 $-14$ LUFS 或 EBU R128 $-23$ LUFS 上交付，GUI 没有响度目标输入，只能走 CLI/MCP；两条预设都带 $-1$ dBTP 上限，目标判定容差 $0.5$ LU，且只接受 $44.1/48/88.2/96$ kHz，其他采样率直接失败，不会静默按 $48$ kHz 算。

\section{导出前检查}

导出前按下面的顺序过一遍。每一步都不需要“再听一遍整首”，但要能指出“这一项为什么可以过”。

\begin{table}[htbp]
\centering
\small
\begin{tabular}{@{}p{0.05\linewidth}p{0.25\linewidth}p{0.62\linewidth}@{}}
\hline
 & 检查项 & 通过条件 \\
\hline
1 & 工程已保存 & 保存一次；导出不会替你保存工程 \\
2 & 独奏/静音 & 所有独奏关闭；该静音的参考轨、废轨已静音 \\
3 & 主总线 & 推子在 $0$ dB；插入链里没有误开的设备 \\
4 & 余量 & 导出前主峰值留在 $-6$ 至 $-3$ dBFS；没有贴 $0$ \\
5 & 格式 & 采样率与工程一致；交付用 24-bit，中间用 32f \\
6 & 时长 & 默认是工程可闻末端；清空会变成 8 秒，别误清空 \\
7 & 限制器 & 只在需要兜底时开；开了也要知道它是固定参数 \\
8 & 实测 & 导出后读真峰值、积分响度、LRA；未测出不能当 0 \\
\hline
\end{tabular}
\caption{导出前检查清单}
\label{tab:walk-export-check}
\end{table}

如果导出后看到“母带削波 N 个样本，建议减 X dB”，先按建议把主总线或源轨降下来再导，不要用导出限制器硬压。导出限制器只保证它自己的天花板附近，不保证“绝对不削顶”；$16$ 倍真峰值仍有固有欠读，交付时留 $1$ dB 左右的余量，不要把读数贴在 $0$。

\section{常见误判与边界}

\begin{table}[htbp]
\centering
\small
\begin{tabular}{@{}p{0.36\linewidth}p{0.58\linewidth}@{}}
\hline
容易误判 & 实际 \\
\hline
“推子就是增益结构” & 推子是电平控制；通道条还有独立的输入/输出增益，源素材的削波不能靠拉推子修好 \\
“主表 LUFS 就是交付积分” & 主表 LUFS 是平滑 RMS 减 $0.691$ dB，不是带门限的积分响度 \\
“压缩器 GR 0.0 dB 就是没压” & 那行 GR 是固定文案，不是实时增益衰减表；用旁通 A/B 判断 \\
“限制器行 $-0.1$ dBTP 就是实际限制” & 该行是静态占位；没有设备时也会显示，不能当真实读数 \\
“EQ 是一个旋钮” & GUI 只暴露低架增益；频率、中频、高架走 MCP/JSON \\
“装混响就能调大小/湿度” & GUI 只能插入/旁通；参数走 MCP/JSON \\
“共享混响用默认值就行” & 默认湿声 $0.25$ 会把 $75\%$ 干声再送一遍，造成双干声；共享返回要把湿声推到接近 $1.0$ \\
“发送有 pre/post” & 发送边只有增益，没有前/后级开关 \\
“GUI 能新建自动化泳道” & 只能点已有泳道；新建/删除/拖点/改曲线走 MCP/JSON \\
“导出限制器就是响度归一化” & 它只是固定参数兜底；响度目标只在 CLI/MCP \\
“真峰值保证绝对不削顶” & $16$ 倍真峰值有固有欠读；留余量，不把读数当数学保证 \\
\hline
\end{tabular}
\caption{从粗混到成品的常见误判}
\label{tab:walk-pitfalls}
\end{table}

这张表也是本章的边界：夜半没有自动母带、自动响度匹配、参考曲电平对齐、发送 pre/post、抖动模式/噪声整形、主总线一键压缩器；没有的能力不要写成“以后会有”，也不要拿别的 DAW 的入口来套。第 9 章的总线链顺序仍然有效：主总线推子在插入链之前，导出限制器在插入链之后；本章只加“怎么走”，不改那条链。

\section{练习}

\begin{enumerate}
  \item 选一首 8 小节，从粗混走到成品；每一步只改一个变量，并在工程备注写“听感—证据—改动—回听”。
  \item 至少做一次“先关限制器导出读基线，再开限制器导出”的对比；不要凭控制室占位读数判断响度。
  \item 交付前做一次完整回听：正常、DIM、MONO、SIDE 四个条件；记录仍未解决的问题，不要静默略过。
\end{enumerate}

\section{术语、截图与帮助映射}

\begin{table}[htbp]
\centering
\small
\begin{tabular}{@{}p{0.22\linewidth}p{0.26\linewidth}p{0.22\linewidth}p{0.26\linewidth}@{}}
\hline
中文 & English & 中文 & English \\
\hline
粗混 & rough mix & 静态平衡 & static balance \\
增益结构 & gain staging & 参考曲 & reference track \\
低架 & low shelf & 中频峰 & mid peak \\
高架 & high shelf & 软膝 & soft knee \\
检波窗口 & detector window & 补偿增益 & makeup gain \\
辅助返回 & auxiliary return & 发送边 & send edge \\
导出限制器 & export limiter & 门限积分响度 & gated integrated loudness \\
真峰值 & true peak & 响度范围 & loudness range (LRA) \\
抖动 & dither & 母线 & master bus \\
\hline
\end{tabular}
\caption{本章新增术语的中英对照}
\label{tab:walk-terms}
\end{table}

\begin{table}[htbp]
\centering
\small
\begin{tabular}{@{}p{0.21\linewidth}p{0.40\linewidth}p{0.35\linewidth}@{}}
\hline
图名 & 界面状态 & 操作步骤 \\
\hline
\texttt{walkthrough-console} & 混音控制台至少六条轨，推子、声像、轨道电平与主总线电平同屏 & 按 X 或 F8 打开混音控制台，播放主歌一，截图 \\
\texttt{walkthrough-eq-low} & 设备机架 EQ 面板只显示低架增益滑块 & 选中主唱轨，装 EQ，拖动低架增益，截图 \\
\texttt{walkthrough-compressor} & 压缩器行显示阈值、比率与 GR 文案 & 轨道菜单插 Compressor，点压缩器行循环比率，截图 \\
\texttt{walkthrough-automation} & 编曲区音量泳道有段落点，主唱轨与副歌边界可见 & 按 A 展开自动化，点泳道加/改点，截图 \\
\texttt{walkthrough-export} & WAV 导出对话框的五个控件与导出后实测行同屏 & 文件菜单导出 WAV，导出一次后截图 \\
\hline
\end{tabular}
\caption{本章截图清单}
\label{tab:walk-screens}
\end{table}

\begin{table}[htbp]
\centering
\small
\begin{tabular}{@{}p{0.34\linewidth}p{0.52\linewidth}@{}}
\hline
本章小节 & 应用内帮助锚点 \\
\hline
准备与参考曲 & \texttt{walkthrough-prep} \\
粗平衡与增益结构 & \texttt{walkthrough-balance-gain} \\
EQ、动态与空间 & \texttt{walkthrough-processing} \\
自动化 & \texttt{walkthrough-automation} \\
总线与导出 & \texttt{walkthrough-bus-export} \\
导出前检查 & \texttt{walkthrough-export-check} \\
常见误判与边界 & \texttt{walkthrough-boundaries} \\
\hline
\end{tabular}
\caption{第十章帮助主题映射}
\label{tab:walk-help}
\end{table}
